\documentclass[12pt]{article}
\usepackage[margin=0.58in,top=0.62in,bottom=0.48in]{geometry}
\usepackage{lmodern}
\usepackage{microtype}
\usepackage{fancyhdr}
\usepackage{titlesec}
\usepackage{enumitem}
\usepackage{lastpage}
\usepackage[hidelinks]{hyperref}

\setlength{\parindent}{0pt}
\setlength{\parskip}{0.18em}
\setlength{\headheight}{26pt}
\raggedbottom
\titleformat{\section}{\normalsize\bfseries\scshape}{}{0pt}{}
\titleformat{\subsection}{\normalsize\bfseries}{}{0pt}{}
\titlespacing*{\section}{0pt}{0.35em}{0.12em}
\titlespacing*{\subsection}{0pt}{0.25em}{0.08em}
\pagestyle{fancy}
\fancyhf{}
\fancyhead[C]{\footnotesize Lit Example Reader \quad Assignment ID: U1L16LE \quad Page \thepage\ of \pageref{LastPage}}
\renewcommand{\headrulewidth}{0.5pt}

\newenvironment{playpassage}{\begin{quote}\footnotesize\setlength{\parskip}{0.08em}\setlength{\parindent}{0pt}}{\end{quote}}
\newcommand{\speaker}[1]{\textbf{#1.}}

\begin{document}

\begin{center}
{\LARGE\bfseries Week 16 Lit Example Reader}\par\vspace{0.02em}
{\large\scshape Shakespeare: Rhetoric and Proof}\par\vspace{0.06em}
\rule{0.78\textwidth}{0.5pt}
\end{center}

\textbf{Purpose:} Shakespeare passages for Week 16: separate what a speech makes an audience feel from what it actually proves.

\textbf{What to watch for:} Speakers use grief, irony, manhood, shame, tokens, and public character. Ask what claim is under proof and whether rhetorical force matches, exceeds, or replaces that proof.

\section*{Before the Passages: The Week 16 Logic Tool}

\begin{itemize}[leftmargin=1.3em,itemsep=0.05em,topsep=0.08em,parsep=0.03em]
\item \textbf{Persuasion:} audience movement---belief, feeling, action.
\item \textbf{Proof:} reasons and evidence that support the claim.
\item \textbf{Logos / ethos / pathos:} reason, character, emotion.
\item \textbf{Rhetoric-vs-proof map:} claim; audience; appeals; proved; only suggested; force vs proof.
\end{itemize}

\section*{Passage 1: Julius Caesar, Act 3, Scene 2 --- Antony at the Funeral}

\textbf{Context:} Brutus has justified Caesar's death. Antony speaks next. He repeatedly calls the conspirators honorable while showing Caesar's mantle, wounds, and will to the crowd.

\begin{playpassage}
\speaker{Antony} Friends, Romans, countrymen, lend me your ears;\\
I come to bury Caesar, not to praise him.\\
...\\
He was my friend, faithful and just to me:\\
But Brutus says he was ambitious;\\
And Brutus is an honourable man.\\
...\\
When that the poor have cried, Caesar hath wept:\\
Ambition should be made of sterner stuff:\\
Yet Brutus says he was ambitious;\\
And Brutus is an honourable man.\\
...\\
O, what a fall was there, my countrymen!\\
Then I, and you, and all of us fell down,\\
Whilst bloody treason flourish'd over us.
\end{playpassage}

\subsection*{Reader Notes}

This is your assisted example. Antony's claim-pressure is not only ``Caesar was not ambitious.'' Watch how ethos (friendship, reluctance to praise), pathos (tears, wounds, shared fall), and selective logos (examples against ambition) work together. The crowd's movement is real. That does not automatically settle whether Caesar had to die, whether Brutus is guilty of bloody treason, or what Rome should do next. Separate force from proof.

\textbf{Reading task:} Underline one ethos move, one pathos move, and one logos-style example. In the margin, write the crowd response Antony seems to want in five words or fewer.

\newpage

\section*{Passage 2: Macbeth, Act 1, Scene 7 --- Lady Macbeth's Pressure}

\textbf{Context:} Macbeth hesitates before murdering Duncan. Lady Macbeth answers with manhood, shame, vivid courage-language, and a practical cover-up plan.

\begin{playpassage}
\speaker{Lady Macbeth} Was the hope drunk\\
Wherein you dress'd yourself? hath it slept since?\\
And wakes it now, to look so green and pale\\
At what it did so freely? From this time\\
Such I account thy love. Art thou afeard\\
To be the same in thine own act and valour\\
As thou art in desire?\\
...\\
When you durst do it, then you were a man;\\
...\\
Screw your courage to the sticking-place,\\
And we'll not fail.
\end{playpassage}

\subsection*{Reader Notes}

This passage shifts more work to you. Do not stop at calling Lady Macbeth ``manipulative.'' Mark what she wants Macbeth to conclude, which appeals she uses, and what moral or factual burden still remains after the speech lands.

\textbf{Reading task:} Box the words that pressure identity or shame. Star one practical promise. Write one question about what still needs proof before murder could be justified.

\section*{How to Use These Passages on the Worksheet}

Do not summarize the speeches as plot. Use Week 16's rhetoric-vs-proof map: main claim, audience/desired response, logos/ethos/pathos, what is proved, what is only suggested, and whether force matches, exceeds, or replaces proof.

\end{document}
